%! AP Statistics — Unit 2, Lesson 2.3 — Group Activity (Answer Key)
\documentclass[10pt]{article}
\usepackage{apstats-article}
\usepackage{apstats-key}

%=============================================================================
\begin{document}

\pageheader{Unit 2, Lesson 2.3}{Group Activity --- Answer Key}
\begin{tabularx}{\linewidth}{X X X}
Name:\;\ans{Key} & Partner:\; & Period:\;
\end{tabularx}

\vspace{0.18in}

\begin{tcolorbox}[colback=sky, colframe=navy, boxrule=0.9pt, arc=2mm,
    left=4mm, right=4mm, top=3mm, bottom=3mm]
\small\textbf{Context:} Sleep and GPA study, $n = 320$.

\vspace{8pt}
\renewcommand{\arraystretch}{1.9}
\begin{tabularx}{\linewidth}{@{} l C{3.2cm} C{3.2cm} C{2.4cm} @{}}
\rowcolor{sky}
                          & \textbf{GPA $\geq$ 3.0} & \textbf{GPA $<$ 3.0} & \textbf{Row Total} \\
\hline
Sleep $<$ 7 hrs  & 40 & 80 & \ans{120} \\
Sleep $\geq$ 7 hrs & 108 & 92 & \ans{200} \\
\hline
Col.\ Total      & \ans{148} & \ans{172} & \ans{320} \\
\end{tabularx}
\end{tcolorbox}

\vspace{0.18in}

{\large\bfseries\color{navy} Part 1 \quad Conditional Proportions}

\vspace{0.1in}

\begin{enumerate}

    \item \ans{Row totals: 120, 200. Column totals: 148, 172. Grand total: 320.}

    \vspace{0.08in}

    \item $\hat{p}_{\text{low sleep}} = \dfrac{\ans{40}}{\ans{120}} = \ans{0.333}$
          \hfill
          $\hat{p}_{\text{adequate}} = \dfrac{\ans{108}}{\ans{200}} = \ans{0.54}$

\end{enumerate}

\vspace{0.22in}

{\large\bfseries\color{navy} Part 2 \quad Difference in Proportions}

\vspace{0.1in}

\begin{enumerate}[resume]

    \item $\hat{p}_{\text{adequate}} - \hat{p}_{\text{low}} = \ans{0.54} - \ans{0.333} = \ans{0.207}$

    \vspace{0.16in}

    \item \ansline{The proportion of students with a GPA $\geq$ 3.0 was about 21 percentage
    points higher among students who slept 7 or more hours than among students who slept
    fewer than 7 hours.}

\end{enumerate}

\vspace{0.22in}

{\large\bfseries\color{navy} Part 3 \quad Relative Risk}

\vspace{0.1in}

\begin{enumerate}[resume]

    \item $\text{RR} = \dfrac{\ans{0.54}}{\ans{0.333}} = \ans{1.62}$

    \vspace{0.16in}

    \item \ansline{Students who slept 7 or more hours were about 1.62 times as likely to have
    a GPA $\geq$ 3.0 as students who slept fewer than 7 hours.}

\end{enumerate}

\vspace{0.22in}

{\large\bfseries\color{navy} Part 4 \quad Comparison and Conclusion}

\vspace{0.1in}

\begin{enumerate}[resume]

    \item \ansline{In this sample, students who got adequate sleep (7$+$ hrs) were substantially
    more likely to have a GPA $\geq$ 3.0 than those who slept fewer than 7 hours: the proportion
    with a high GPA was 21 percentage points higher (54\% vs.\ 33\%), and adequate-sleep students
    were 1.62 times as likely to have a high GPA.  Both statistics suggest a meaningful positive
    association between sleep and academic performance.}

    \vspace{0.1in}

    \item \ansline{Both statistics are useful. The difference in proportions (21 pp) is easy
    to interpret on the percentage-point scale. The relative risk (1.62$\times$) may be preferred
    in public-health reporting because it shows the multiplicative effect. Since the outcome
    (GPA $\geq$ 3.0) is fairly common ($\approx$46\%), the absolute difference is also quite
    interpretable, so either is reasonable to emphasize.}

\end{enumerate}

\vspace{0.18in}

\begin{tcolorbox}[colback=redbg, colframe=redacc, boxrule=0.8pt, arc=2mm,
    left=4mm, right=4mm, top=3mm, bottom=3mm]
\small\textbf{AP-Style Practice}

\begin{enumerate}[label=\textbf{(\Alph*)}, leftmargin=2em, itemsep=6pt]
  \item There is no association because neither group had a majority with GPA $\geq$ 3.0.
  \item Students with adequate sleep are more likely to have a GPA $\geq$ 3.0, but the
        relative risk is close to 1 so the association is weak.
  \item Students with adequate sleep are more likely to have a GPA $\geq$ 3.0; both the
        difference in proportions and relative risk indicate a meaningful association.
        \textcolor{keyred}{\textbf{$\leftarrow$ correct}}
  \item The relative risk cannot be interpreted without knowing the marginal proportions.
\end{enumerate}

\vspace{6pt}
\ans{(C). The difference in proportions ($\approx$0.21) is large, and RR $= 1.62$ is
meaningfully above 1. (B) is wrong: RR $= 1.62$ is not close to 1. (D) is wrong: conditional
proportions are sufficient to compute and interpret RR.}
\end{tcolorbox}

\begin{teachernote}
\textbf{Differentiation tiers:}
\begin{itemize}[leftmargin=1.4em, itemsep=2pt]
  \item \textit{Support (Tier R):} Provide the interpretation templates from the guided notes;
        students fill in the numbers and context words.
  \item \textit{Extension (Tier E):} Ask groups to reverse the comparison (low-sleep vs.\
        adequate-sleep as baseline) and verify that RR becomes $1/1.62 \approx 0.62$.
        Discuss: same association, different framing.
\end{itemize}

\textbf{Common error in Part 2:} Students compute $0.333 - 0.54$ (negative result). Accept this
if direction is stated correctly in the interpretation (``low-sleep students are 21 pp \emph{less}
likely$\ldots$''). Emphasize that which group is Group 1 is a choice --- be consistent and state direction.

\textbf{AP item (C):} (B) is a common wrong answer --- a relative risk of 1.62 is not close to 1.
Use to reinforce that ``close to 1'' means roughly 1.00--1.05 in typical contexts.
\end{teachernote}

\end{document}
